\section{Additional details from the Distance Network: calibration parameters and host/calibrator distances}\label{app:sec:distance_table}

The main result of the {\DN} analysis is of course the value and uncertainty of {\Hcst}, reported in Table~\ref{tab:variants} for all variants, together with the value of the total $ \chi^2 $ to evaluate the quality of the fit.  In addition, the {\DN} process can provide more detailed information on each class of calibrators and on host and the calibrated distances.

\subsection{Calibration parameters}

The {\DN} approach can include several secondary distance indicators whose calibration is determined simultaneously with all other parameters of the solution; each can be used to constrain {\Hcst} via tracers in the Hubble Flow.  For each variant, Table~\ref{tab:extrapars} provides the resulting calibration parameter, the Hubble Flow intercept, and the number of calibrators for each of the main secondary distance indicators: SNe~Ia, SBF, SNe~II, and TF. For SNe~Ia and SN~II, the calibration parameter is generally the reference absolute magnitude of a standardized SN in the appropriate passband; the exception is the BayesSN method, which uses a different definition.  For SBF and TF, the calibration parameter is reported as the \textit{difference} with respect to the published value, as described in Section~\ref{sec:data_appendix}. In all cases, the reported uncertainty is marginalized over all other parameters in the {\DN}.

The Hubble Flow intercept is the parameter $ a_{\itype} $ defined in Equation~\ref{eq:hubble_flow} and discussed in Section~\ref{app:sec:eq_hf}.  It is expressed in units of $ \log(\Hcst) $ (numerically equivalent to 0.2 mag); the reported uncertainty includes the intrinsic and measurement error for the tracers, but does not include the calibration uncertainty.  For some versions of SNe~Ia, the reported uncertainty includes covariance terms between different tracers, as reported in the original sources.

\begin{longrotatetable}
\movetabledown=2cm
\begin{deluxetable}{lccccccccrrrr}
\tabletypesize{\tiny}
\tablecaption {Additional fitted parameters for all variants\label{tab:extrapars}}
\tablehead{
\colhead {Variant}  & \multicolumn{4}{c}{Calibration parameter (mag)\tablenotemark{\textup{(}a\textup{)}}}                      &  \multicolumn{4}{c}{Hubble flow intercept\tablenotemark{\textup{(}b\textup{)}}}                    & \multicolumn{4}{c}{Number of calibrators} \\
\colhead{} &  
\multicolumn{1}{c}{SN~Ia}&\multicolumn{1}{c}{SBF}&\multicolumn{1}{c}{SN~II}&\multicolumn{1}{c}{TF}& 
\multicolumn{1}{c}{SN~Ia}&\multicolumn{1}{c}{SBF}&\multicolumn{1}{c}{SN~II}&\multicolumn{1}{c}{TF}& 
\multicolumn{1}{c}{SN~Ia}&\multicolumn{1}{c}{SBF}&\multicolumn{1}{c}{SN~II}&\multicolumn{1}{c}{TF}
}
\startdata
 V00 & $ -19.252 \pm   0.022 $  & $  -0.023 \pm   0.027 $  &          ---             &          ---             & $   0.716 \pm   0.002 $  & $  -3.128 \pm   0.003 $  &          ---             &          ---             &   55 &   14 &    0 &    0 \\
 V01 & $ -19.252 \pm   0.022 $  & $  -0.024 \pm   0.027 $  &          ---             &          ---             & $   0.716 \pm   0.002 $  & $  -3.128 \pm   0.003 $  &          ---             &          ---             &   55 &   14 &    0 &    0 \\
 V02 & $ -19.252 \pm   0.022 $  & $  -0.023 \pm   0.027 $  &          ---             &          ---             & $   0.716 \pm   0.002 $  & $  -3.128 \pm   0.003 $  &          ---             &          ---             &   55 &   14 &    0 &    0 \\
 V03 & $ -19.252 \pm   0.022 $  & $  -0.021 \pm   0.027 $  &          ---             &          ---             & $   0.716 \pm   0.002 $  & $  -3.128 \pm   0.003 $  &          ---             &          ---             &   55 &   14 &    0 &    0 \\
 V04 & $ -19.251 \pm   0.022 $  & $  -0.023 \pm   0.027 $  & $ -16.760 \pm   0.044 $  &          ---             & $   0.716 \pm   0.002 $  & $  -3.128 \pm   0.003 $  & $   0.219 \pm   0.008 $  &          ---             &   55 &   14 &    9 &    0 \\
 V05 & $ -19.246 \pm   0.021 $  & $  -0.017 \pm   0.025 $  &          ---             &          ---             & $   0.716 \pm   0.002 $  & $  -3.128 \pm   0.003 $  &          ---             &          ---             &   55 &   14 &    0 &    0 \\
 V06 & $ -19.241 \pm   0.022 $  & $  -0.012 \pm   0.026 $  &          ---             & $  -0.076 \pm   0.024 $  & $   0.716 \pm   0.002 $  & $  -3.128 \pm   0.003 $  &          ---             & $  -3.116 \pm   0.001 $  &   55 &   14 &    0 &   74 \\
 V07 & $ -19.257 \pm   0.022 $  & $  -0.029 \pm   0.027 $  &          ---             &          ---             & $   0.716 \pm   0.002 $  & $  -3.128 \pm   0.003 $  &          ---             &          ---             &   55 &   14 &    0 &    0 \\
 V08 & $ -19.282 \pm   0.038 $  & $  -0.051 \pm   0.040 $  &          ---             &          ---             & $   0.716 \pm   0.002 $  & $  -3.128 \pm   0.003 $  &          ---             &          ---             &   35 &   14 &    0 &    0 \\
V08B & $ -19.271 \pm   0.025 $  & $  -0.041 \pm   0.029 $  &          ---             &          ---             & $   0.716 \pm   0.002 $  & $  -3.128 \pm   0.003 $  &          ---             &          ---             &   35 &   14 &    0 &    0 \\
 V09 & $ -19.249 \pm   0.024 $  &          ---             &          ---             &          ---             & $   0.716 \pm   0.002 $  &          ---             &          ---             &          ---             &   42 &    0 &    0 &    0 \\
V09B & $ -19.243 \pm   0.023 $  & $  -0.020 \pm   0.028 $  &          ---             &          ---             & $   0.716 \pm   0.002 $  & $  -3.128 \pm   0.003 $  &          ---             &          ---             &   42 &   14 &    0 &    0 \\
 V10 & $ -19.249 \pm   0.026 $  & $  -0.020 \pm   0.030 $  &          ---             &          ---             & $   0.716 \pm   0.002 $  & $  -3.128 \pm   0.003 $  &          ---             &          ---             &   55 &   14 &    0 &    0 \\
 V11 & $ -19.257 \pm   0.024 $  & $  -0.028 \pm   0.028 $  &          ---             &          ---             & $   0.716 \pm   0.002 $  & $  -3.128 \pm   0.003 $  &          ---             &          ---             &   55 &   14 &    0 &    0 \\
 V12 & $ -19.262 \pm   0.026 $  &          ---             &          ---             &          ---             & $   0.716 \pm   0.002 $  &          ---             &          ---             &          ---             &   42 &    0 &    0 &    0 \\
V12B & $ -19.246 \pm   0.024 $  &          ---             &          ---             &          ---             & $   0.716 \pm   0.002 $  &          ---             &          ---             &          ---             &   42 &    0 &    0 &    0 \\
 V13 &          ---             & $  -0.029 \pm   0.052 $  &          ---             &          ---             &          ---             & $  -3.128 \pm   0.003 $  &          ---             &          ---             &    0 &   14 &    0 &    0 \\
 V14 & $ -19.255 \pm   0.023 $  &          ---             &          ---             &          ---             & $   0.716 \pm   0.002 $  &          ---             &          ---             &          ---             &   55 &    0 &    0 &    0 \\
 V15 & $ -19.248 \pm   0.023 $  & $  -0.020 \pm   0.028 $  &          ---             &          ---             & $   0.716 \pm   0.002 $  & $  -3.128 \pm   0.003 $  &          ---             &          ---             &   55 &   14 &    0 &    0 \\
 V16 & $ -19.244 \pm   0.038 $  & $  -0.017 \pm   0.040 $  &          ---             &          ---             & $   0.716 \pm   0.002 $  & $  -3.128 \pm   0.003 $  &          ---             &          ---             &   24 &   14 &    0 &    0 \\
V16B & $ -19.236 \pm   0.026 $  & $  -0.009 \pm   0.030 $  &          ---             &          ---             & $   0.716 \pm   0.002 $  & $  -3.128 \pm   0.003 $  &          ---             &          ---             &   24 &   14 &    0 &    0 \\
 V17 & $ -19.267 \pm   0.024 $  & $  -0.039 \pm   0.029 $  &          ---             &          ---             & $   0.716 \pm   0.002 $  & $  -3.128 \pm   0.003 $  &          ---             &          ---             &   55 &   13 &    0 &    0 \\
 V18 & $ -19.246 \pm   0.023 $  & $  -0.019 \pm   0.027 $  &          ---             &          ---             & $   0.716 \pm   0.002 $  & $  -3.128 \pm   0.003 $  &          ---             &          ---             &   48 &   14 &    0 &    0 \\
 V19 & $ -19.249 \pm   0.022 $  & $  -0.027 \pm   0.027 $  &          ---             &          ---             & $   0.713 \pm   0.002 $  & $  -3.130 \pm   0.003 $  &          ---             &          ---             &   55 &   14 &    0 &    0 \\
 V20 & $ -19.255 \pm   0.023 $  &          ---             &          ---             &          ---             & $   0.716 \pm   0.002 $  &          ---             &          ---             &          ---             &   55 &    0 &    0 &    0 \\
 V21 & $ -19.254 \pm   0.023 $  &          ---             &          ---             &          ---             & $   0.712 \pm   0.003 $  &          ---             &          ---             &          ---             &   55 &    0 &    0 &    0 \\
 V22 & $ -19.122 \pm   0.027 $  & $  -0.030 \pm   0.031 $  &          ---             &          ---             & $   0.689 \pm   0.003 $  & $  -3.128 \pm   0.003 $  &          ---             &          ---             &   55 &   14 &    0 &    0 \\
 V23\tablenotemark{\textup{(}c\textup{)}} &
       $   0.014 \pm   0.029 $ & $  -0.017 \pm   0.033 $  &          ---             &          ---             & $  -3.135 \pm   0.004 $  & $  -3.128 \pm   0.003 $  &          ---             &          ---             &   50 &   14 &    0 &    0 \\
 V24 & $ -19.293 \pm   0.022 $  & $  -0.024 \pm   0.027 $  &          ---             &          ---             & $   0.724 \pm   0.002 $  & $  -3.128 \pm   0.003 $  &          ---             &          ---             &   55 &   14 &    0 &    0 \\
 V25 & $ -18.354 \pm   0.030 $  & $  -0.055 \pm   0.034 $  &          ---             &          ---             & $   0.529 \pm   0.004 $  & $  -3.128 \pm   0.003 $  &          ---             &          ---             &   14 &   14 &    0 &    0 \\
 V26 & $ -18.567 \pm   0.033 $  & $  -0.048 \pm   0.036 $  &          ---             &          ---             & $   0.574 \pm   0.004 $  & $  -3.128 \pm   0.003 $  &          ---             &          ---             &   17 &   14 &    0 &    0 \\
 V27 & $ -19.251 \pm   0.022 $  & $  -0.026 \pm   0.027 $  &          ---             &          ---             & $   0.716 \pm   0.002 $  & $  -3.128 \pm   0.003 $  &          ---             &          ---             &   55 &   14 &    0 &    0 \\
 V28 & $ -19.264 \pm   0.022 $  & $  -0.036 \pm   0.027 $  &          ---             &          ---             & $   0.716 \pm   0.002 $  & $  -3.128 \pm   0.003 $  &          ---             &          ---             &   55 &   14 &    0 &    0 \\
 V99 & $ -19.243 \pm   0.020 $  & $  -0.011 \pm   0.024 $  & $ -16.749 \pm   0.042 $  & $  -0.075 \pm   0.021 $  & $   0.716 \pm   0.002 $  & $  -3.128 \pm   0.003 $  & $   0.219 \pm   0.008 $  & $  -3.116 \pm   0.001 $  &   55 &   14 &    9 &   74 \\
V99A & $ -19.251 \pm   0.020 $  & $  -0.019 \pm   0.025 $  & $ -16.757 \pm   0.042 $  &          ---             & $   0.716 \pm   0.002 $  & $  -3.128 \pm   0.003 $  & $   0.219 \pm   0.008 $  &          ---             &   55 &   14 &    9 &    0 \\
\enddata
\tablenotetext{\textup{(}a\textup{)}}{\quad Absolute magnitude for SNe, offset to the nominal calibration for SBF and TF. Uncertainty is marginalized over all other parameters.}
\tablenotetext{\textup{(}b\textup{)}}{\quad In units of $\log_{10}(\Hcst)$. Uncertainty includes only Hubble Flow terms, and does not include the calibration uncertainty.}
\tablenotetext{\textup{(}c\textup{)}}{\quad Calibration and intercept for SNe~Ia measured with the BayesSN method have a different, internally consistent definition.}
\end{deluxetable}
\end{longrotatetable}

\subsection{Host and calibrator distances}

One of the secondary results of the {\DN} approach is an optimized set of distance moduli for all hosts and calibrators included in the Network.  These distances are based not only on the multiple measurements included in the {\DN}, but also on the consistency equations for the various classes of calibrators to which they belong.  They represent the values of the distances that, together, produce the best-fitting solution to all of the {\DN} equations.  As a result, the estimated distances for each host will differ slightly depending on which classes of calibrators are included in the solution, even when no additional measurements are introduced.

Table~\ref{tab:host_distances} reports the best distance estimate for each of the hosts (or host/calibrators) included in two solutions, Baseline (V00) and ``Everything,'' the solution including all methods (V99).  Since hosts not tied to any calibrators are excluded from each solution, many of the hosts do not have a distance estimate in the Baseline version.  SN calibrators are not included, since their distances are not separate parameters; they are assumed identical to the distances to their host, which as a consequence are also affected by the SN calibration via the (one or more) SNe they host.  On the other hand, SBF calibrator galaxies in Fornax and Virgo \textit{are} included individually; their distances are allowed to deviate from the group distance, because of depth effects, and are affected by their SBF luminosity measurements via their global calibration.  The quoted distance uncertainties are marginalized over all parameters of the respective {\DN} solution, and therefore are \textit{not} independent.

\startlongtable
\begin{deluxetable}{lcccc}
\tabletypesize{\scriptsize}
\tablecaption{Distance moduli and statistical uncertainties for host and calibrator systems in the {\DN}.\label{tab:host_distances}}
\tablehead{Name & \multicolumn{2}{c}{Baseline solution} & \multicolumn{2}{c}{``Everything'' solution} \\
           & Value & Error & Value & Error}
\startdata
       IC 2006     &   31.469 &    0.048     &   31.457 &    0.046 \\
          M 33     &    ---   &    ---       &   24.847 &    0.075 \\
          M 64     &    ---   &    ---       &   28.219 &    0.042 \\
          M 66     &   30.200 &    0.045     &   30.202 &    0.044 \\
          M 81     &    ---   &    ---       &   27.901 &    0.081 \\
          M 96     &   30.287 &    0.045     &   30.289 &    0.044 \\
         M 101     &   29.120 &    0.026     &   29.125 &    0.023 \\
         M 104     &    ---   &    ---       &   29.851 &    0.041 \\
      Mrk 1337     &   32.930 &    0.082     &   32.922 &    0.082 \\
        NGC 24     &    ---   &    ---       &   29.315 &    0.036 \\
       NGC 105     &   34.507 &    0.119     &   34.501 &    0.119 \\
       NGC 247     &    ---   &    ---       &   27.854 &    0.036 \\
       NGC 253     &    ---   &    ---       &   27.844 &    0.036 \\
       NGC 300     &    ---   &    ---       &   26.592 &    0.058 \\
       NGC 628     &    ---   &    ---       &   29.897 &    0.122 \\
       NGC 672     &    ---   &    ---       &   29.288 &    0.084 \\
       NGC 691     &   32.836 &    0.089     &   32.821 &    0.085 \\
       NGC 891     &    ---   &    ---       &   30.005 &    0.074 \\
       NGC 976     &   33.566 &    0.087     &   33.553 &    0.085 \\
      NGC 1015     &   32.600 &    0.062     &   32.593 &    0.062 \\
      NGC 1309     &   32.498 &    0.044     &   32.498 &    0.043 \\
      NGC 1316     &   31.385 &    0.042     &   31.386 &    0.041 \\
      NGC 1344     &   31.439 &    0.049     &   31.427 &    0.047 \\
      NGC 1365     &   31.335 &    0.035     &   31.332 &    0.029 \\
      NGC 1374     &   31.479 &    0.049     &   31.466 &    0.048 \\
      NGC 1375     &   31.457 &    0.052     &   31.445 &    0.051 \\
      NGC 1380     &   31.417 &    0.037     &   31.406 &    0.036 \\
      NGC 1399     &   31.504 &    0.039     &   31.493 &    0.037 \\
      NGC 1404     &   31.407 &    0.037     &   31.396 &    0.035 \\
      NGC 1448     &   31.308 &    0.028     &   31.313 &    0.025 \\
      NGC 1559     &   31.423 &    0.030     &   31.422 &    0.026 \\
      NGC 1560     &    ---   &    ---       &   27.539 &    0.148 \\
      NGC 2188     &    ---   &    ---       &   29.618 &    0.050 \\
      NGC 2403     &    ---   &    ---       &   27.528 &    0.058 \\
      NGC 2442     &   31.499 &    0.049     &   31.530 &    0.036 \\
      NGC 2525     &   31.929 &    0.052     &   31.912 &    0.051 \\
      NGC 2608     &   32.648 &    0.115     &   32.560 &    0.111 \\
      NGC 2683     &    ---   &    ---       &   29.979 &    0.057 \\
      NGC 2903     &    ---   &    ---       &   29.817 &    0.049 \\
      NGC 2915     &    ---   &    ---       &   28.191 &    0.076 \\
      NGC 2976     &    ---   &    ---       &   27.822 &    0.050 \\
      NGC 3021     &   32.242 &    0.049     &   32.237 &    0.048 \\
      NGC 3109     &    ---   &    ---       &   25.712 &    0.094 \\
      NGC 3147     &   33.094 &    0.084     &   33.086 &    0.083 \\
      NGC 3254     &   32.396 &    0.057     &   32.410 &    0.055 \\
      NGC 3368     &    ---   &    ---       &   30.245 &    0.036 \\
      NGC 3370     &   32.189 &    0.036     &   32.193 &    0.035 \\
      NGC 3432     &    ---   &    ---       &   29.929 &    0.161 \\
      NGC 3447     &   31.956 &    0.029     &   31.951 &    0.027 \\
      NGC 3583     &   32.783 &    0.064     &   32.782 &    0.063 \\
      NGC 3621     &    ---   &    ---       &   29.224 &    0.042 \\
      NGC 3627     &    ---   &    ---       &   30.259 &    0.042 \\
      NGC 3972     &   31.718 &    0.051     &   31.716 &    0.050 \\
      NGC 3982     &   31.600 &    0.053     &   31.599 &    0.052 \\
      NGC 4038     &   31.672 &    0.030     &   31.673 &    0.029 \\
      NGC 4144     &    ---   &    ---       &   29.194 &    0.042 \\
      NGC 4151     &    ---   &    ---       &   31.036 &    0.070 \\
      NGC 4236     &    ---   &    ---       &   28.298 &    0.050 \\
      NGC 4242     &    ---   &    ---       &   29.403 &    0.076 \\
      NGC 4244     &    ---   &    ---       &   28.184 &    0.075 \\
      NGC 4298     &    ---   &    ---       &   30.842 &    0.167 \\
      NGC 4414     &   31.264 &    0.081     &   31.249 &    0.078 \\
      NGC 4424     &   30.910 &    0.034     &   30.909 &    0.033 \\
      NGC 4455     &    ---   &    ---       &   29.120 &    0.076 \\
      NGC 4457     &   31.076 &    0.086     &   31.074 &    0.086 \\
      NGC 4458     &   31.097 &    0.056     &   31.085 &    0.055 \\
      NGC 4472     &   31.088 &    0.039     &   31.078 &    0.037 \\
      NGC 4489     &   31.002 &    0.056     &   30.990 &    0.055 \\
      NGC 4517     &    ---   &    ---       &   29.663 &    0.135 \\
      NGC 4526     &   30.936 &    0.066     &   30.937 &    0.066 \\
      NGC 4536     &   30.892 &    0.034     &   30.906 &    0.028 \\
      NGC 4552     &   30.978 &    0.041     &   30.968 &    0.040 \\
      NGC 4559     &    ---   &    ---       &   29.733 &    0.042 \\
      NGC 4565     &    ---   &    ---       &   30.415 &    0.036 \\
      NGC 4592     &    ---   &    ---       &   29.749 &    0.050 \\
      NGC 4605     &    ---   &    ---       &   28.723 &    0.036 \\
      NGC 4631     &    ---   &    ---       &   29.333 &    0.036 \\
      NGC 4636     &   31.111 &    0.039     &   31.100 &    0.038 \\
      NGC 4639     &   31.772 &    0.048     &   31.728 &    0.035 \\
      NGC 4649     &   31.071 &    0.040     &   31.060 &    0.038 \\
      NGC 4666     &   30.903 &    0.057     &   30.903 &    0.057 \\
      NGC 4680     &   32.510 &    0.146     &   32.508 &    0.146 \\
      NGC 4697     &   30.347 &    0.041     &   30.337 &    0.039 \\
      NGC 4945     &    ---   &    ---       &   27.765 &    0.058 \\
      NGC 5023     &    ---   &    ---       &   28.926 &    0.058 \\
      NGC 5055     &    ---   &    ---       &   29.773 &    0.036 \\
      NGC 5194     &    ---   &    ---       &   29.619 &    0.059 \\
      NGC 5204     &    ---   &    ---       &   28.443 &    0.216 \\
      NGC 5398     &    ---   &    ---       &   30.320 &    0.185 \\
      NGC 5468     &   33.050 &    0.031     &   33.050 &    0.030 \\
      NGC 5584     &   31.853 &    0.030     &   31.861 &    0.026 \\
      NGC 5643     &   30.535 &    0.027     &   30.533 &    0.023 \\
      NGC 5728     &   32.842 &    0.102     &   32.819 &    0.094 \\
      NGC 5861     &   32.171 &    0.064     &   32.176 &    0.062 \\
      NGC 5907     &    ---   &    ---       &   31.179 &    0.057 \\
      NGC 5917     &   32.343 &    0.077     &   32.339 &    0.077 \\
      NGC 6503     &    ---   &    ---       &   28.997 &    0.066 \\
      NGC 6744     &    ---   &    ---       &   29.861 &    0.066 \\
      NGC 6946     &    ---   &    ---       &   29.206 &    0.089 \\
      NGC 7090     &    ---   &    ---       &   29.887 &    0.058 \\
      NGC 7250     &   31.603 &    0.048     &   31.570 &    0.036 \\
      NGC 7329     &   33.269 &    0.069     &   33.264 &    0.068 \\
      NGC 7541     &   32.579 &    0.089     &   32.577 &    0.089 \\
      NGC 7640     &    ---   &    ---       &   29.652 &    0.066 \\
      NGC 7678     &   33.300 &    0.082     &   33.295 &    0.082 \\
      NGC 7793     &    ---   &    ---       &   27.788 &    0.053 \\
      NGC 7814     &   30.926 &    0.071     &   30.948 &    0.068 \\
      PGC 6574     &    ---   &    ---       &   29.360 &    0.076 \\
      PGC 9962     &    ---   &    ---       &   28.535 &    0.036 \\
     PGC 11139     &    ---   &    ---       &   28.844 &    0.036 \\
     PGC 13163     &    ---   &    ---       &   28.966 &    0.042 \\
     PGC 16957     &    ---   &    ---       &   27.891 &    0.199 \\
     PGC 45084     &    ---   &    ---       &   28.937 &    0.067 \\
     PGC 47847     &    ---   &    ---       &   29.211 &    0.050 \\
     PGC 54392     &    ---   &    ---       &   27.283 &    0.103 \\
     PGC 64181     &    ---   &    ---       &   29.838 &    0.122 \\
     PGC 65603     &    ---   &    ---       &   28.699 &    0.058 \\
      UGC 1281     &    ---   &    ---       &   28.615 &    0.036 \\
      UGC 9391     &   32.818 &    0.056     &   32.817 &    0.056 \\
      UGCA 319     &    ---   &    ---       &   28.853 &    0.050 \\

\enddata
\end{deluxetable}
